\section{Discussion}

The current evidence supports a focused interpretation of \method{}. A fixed-size anchored scene representation can preserve point-feature structure while allowing query images to reuse the same cached memory. The Indoor6 DINOv2+MapAnything results and the ACE-FCN/GLACE-compatible experiments test this representation under different feature sources and coordinate-regression heads; neither branch changes the central token interface.

The main limitation is evidence coverage rather than a change in method scope. The final paper must consolidate all-scene comparisons under the fixed global protocol, report the capacity and efficiency trade-off of the 64-token representation, and support any cross-view-consistency claim with a dedicated statistic rather than qualitative examples alone.

Several limitations remain. The Indoor6 ACE-FCN+GLACE evidence requires a clean all-scene summary before final submission. Wayspots full-scene aggregation is needed to avoid overemphasizing a small subset of scenes. Cambridge and APR+LMC experiments are important for testing outdoor scalability and non-SCR generality, respectively. Finally, the relationship between \method{} and ACE-G map-code representations should be evaluated directly or clearly described as future work.
